% Appendix B — polished, paste-ready replacement for chapters/93-AppendixB.tex.
% Parameter values read from Neuromechanical_Models/Biped_2xCPG_wSubs walker.

\chapter{Neuron Model Equations}\label{app:neuron_equations}

This appendix documents the neuron and synapse models used by the synthetic
nervous systems in this project, both in the current AnimatLab walker models and
as they will be instantiated in SNS-Toolbox for the MuJoCo pipeline of
Chapter~\ref{ch:futurework}.

\section{Spiking Neuron Model}

Each spiking neuron is a conductance-based leaky integrator with threshold and
adaptation,
\begin{equation}
    C_{m}\,\dot{V} \;=\; -\,g_{\text{leak}}\left(V - E_{\text{rest}}\right)
                      \;+\; I_{\text{syn}} \;+\; I_{\text{app}},
    \label{eq:app_membrane}
\end{equation}
where $V$ is the membrane potential, $C_{m}$ the membrane capacitance,
$g_{\text{leak}} = C_{m}/\tau_{m}$ the leak conductance set by the membrane time
constant $\tau_{m}$, $E_{\text{rest}}$ the resting potential, $I_{\text{syn}}$
the summed synaptic current, and $I_{\text{app}}$ any applied tonic stimulation
current. A spike is emitted when $V$ crosses the firing threshold, after which
$V$ resets and an after-hyperpolarization of amplitude $A_{\text{AHP}}$ with
time constant $\tau_{\text{AHP}}$ is applied. The threshold accommodates with
firing history (relative accommodation $r_{\text{acc}}$, time constant
$\tau_{\text{acc}}$).

\begin{table}[htbp]
\caption{Spiking neuron parameters for the rhythm-generator,
pattern-formation, and motoneuron populations of the AnimatLab bipedal walker
model. The sensory afferent (MV) populations differ, as noted in the text.}
\label{tab:app_neuron}
\begin{tabular}{@{}lll@{}}
\toprule
\textbf{Parameter} & \textbf{Symbol} & \textbf{Value} \\
\midrule
Resting potential        & $E_{\text{rest}}$   & \qty{-60}{\mV} \\
Membrane time constant   & $\tau_{m}$          & \qty{5}{\ms}   \\
Initial firing threshold & $V_{\text{th}}$     & \qty{-55}{\mV} (motor neurons: \qty{-59}{\mV}) \\
Relative accommodation   & $r_{\text{acc}}$    & 0.3            \\
Accommodation time constant & $\tau_{\text{acc}}$ & \qty{10}{\ms} \\
AHP amplitude            & $A_{\text{AHP}}$    & \qty{1}{\mV}   \\
AHP time constant        & $\tau_{\text{AHP}}$ & \qty{3}{\ms}   \\
\bottomrule
\end{tabular}
\end{table}

The sensory afferent (MV) populations that carry position and load signals into
the circuit differ from Table~\ref{tab:app_neuron}: they rest at
\qty{-100}{\mV} with a firing threshold of \qty{+50}{\mV}, use no threshold
accommodation, and have a slower \qty{90}{\ms} after-hyperpolarization. A
subset of interneurons likewise disables accommodation
($r_{\text{acc}} = 0$). Optional calcium channels are defined in all neurons
but disabled in the current configuration ($G_{\text{max,Ca}} = 0$).

\section{Non-Spiking Chemical Synapses}

Pattern-formation and motoneuron-drive connections use graded chemical
synapses, whose conductance is a clipped linear ramp of the presynaptic
potential,
\begin{equation}
    g_{\text{syn}}(V_{\text{pre}})
    \;=\; g_{\text{max}}\;
    \operatorname{clamp}\!\left(
        \frac{V_{\text{pre}} - V_{\text{th}}}{V_{\text{sat}} - V_{\text{th}}},
        \,0,\,1\right),
    \qquad
    I_{\text{syn}} = g_{\text{syn}}\left(V_{\text{post}} - E_{\text{syn}}\right),
    \label{eq:app_nonsyn}
\end{equation}
with threshold and saturation potentials $V_{\text{th}}$ and $V_{\text{sat}}$
and reversal potential $E_{\text{syn}}$. The walker model instantiates three
non-spiking synapse types: excitatory nicotinic ACh
($E_{\text{syn}} = \qty{-10}{\mV}$, $g_{\text{max}} = \qty{0.5}{\micro\siemens}$,
$V_{\text{th}} \to V_{\text{sat}} = \qty{-65}{\mV} \to \qty{-20}{\mV}$),
hyperpolarizing IPSP ($E_{\text{syn}} = \qty{-70}{\mV}$,
$\qty{-55}{\mV} \to \qty{-40}{\mV}$), and depolarizing IPSP
($E_{\text{syn}} = \qty{-55}{\mV}$,
$\qty{-60}{\mV} \to \qty{-40}{\mV}$). Spiking chemical synapses and electrical
synapses (rectifying coupling 0.1 in the low state and 0.3 in the high state;
non-rectifying 0.2) are also defined and remain available for the sensory
afferent pathways of the MuJoCo/SNS implementation.

\section{Relation to the SNS-Toolbox Implementation}

These are the standard conductance-based models; SNS-Toolbox
\citep{nourse_sns_2023} implements the same non-spiking and spiking classes
with conductance-based synapses, so the parameter set above transfers directly
to the rhythm-generator, pattern-formation, and feedback circuits planned for
Chapter~\ref{ch:futurework}. The tonic stimulation tests of
Section~\ref{sec:fw_cpg} correspond to constant $I_{\text{app}}$ injections
into Equation~(\ref{eq:app_membrane}) at the motoneuron, pattern-formation, and
rhythm-generator layers respectively.
